\section{Проектирование и разработка веб-интерфейса}

Как уже было сказано ранее, сервис реализует принцип Server-Side Render.
Для работы с полученными сервером данными используются специальные операнды и операции, заключенные
в двойные фигурные скобки.

\begin{listing}[h]
    \htmlfile[firstline=1, lastline=7, firstnumber=1]{../index.html}
    \caption{Начало HTML-документа}
    \label{lst:html-file-beginning}
\end{listing}

В начале документа указываются (листинг \ref{lst:html-file-beginning}):
\begin{enumerate}
    \item Тип документа --- HTML;
    \item Язык документа --- русский;
    \item Технические параметры для браузера:
    \begin{enumerate}
        \item Кодировка UTF-8 для поддержки кириллицы;
        \item Заголовок страницы;
        \item Подключается файл стилей pico.classes.min.css.
    \end{enumerate}
\end{enumerate}

\begin{listing}[h]
    \htmlfile[firstline=9, lastline=13, firstnumber=9]{../index.html}
    \caption{Заголовки страницы}
    \label{lst:html-header}
\end{listing}

В начале тела контейнера main указываются заголовки страницы (листинг \ref{lst:html-header}):
\begin{enumerate}
    \item Дисциплина и номер работы;
    \item Группа и ФИО студента.
\end{enumerate}

\begin{listing}[h]
    \htmlfile[firstline=15, lastline=32, firstnumber=15]{../index.html}
    \caption{Первый артикль}
    \label{lst:html-article1}
\end{listing}

В первом артикле документа создаются кнопки, используемые для вызова SQL-запросов (листинг \ref{lst:html-article1}):
\begin{enumerate}
    \item Строки \texttt{<form action="/query" method="get">} определяют путь запроса и его метод;
    \item Строки \texttt{<input type="hidden" name="opId"};
    \texttt{value="<идентификатор\_операции>">} создают query-параметр \texttt{opId} и присваивают ему значение равное идентификатору операции;
    \item Объекты \texttt{<button type="submit">} создают кнопки и надписи на них.
\end{enumerate}

\begin{listing}[h]
    \htmlfile[firstline=34, lastline=67, firstnumber=34]{../index.html}
    \caption{Логика формирования таблицы}
    \label{lst:html-article2}
\end{listing}

Во втором артикле реализована логика формирования таблицы SQL-запросов (листинг \ref{lst:html-article2}).

Используются элементы \texttt{<table>}, \texttt{<thead>} для инициализации таблицы и заголовков её столбцов.

Далее в двойных фигурных скобках проверяется условие: если поле типа структуры, полученной в результате выполнения,
равен единице, то при помощи элементов \texttt{<th>} создаются заголовки:
фамилия, имя, отчество, номер телефона, зарплата.

В ином случае, если тип равен двум:
фамилия, имя, отчество, адрес.

В ином случае, если тип равен трем:
фамилия, имя, отчество, дата найма.

В ином случае вместо заголовков пишется строка <<Нажмите кнопку>>.

Далее срабатывает цикл \texttt{\{\{ range .Rows \}\}}, который используется для перебора всех элементов массива Rows.
Для каждого элемента создается блок \texttt{<tr>}, инициализирующий строку таблицы.
Внутри каждого блока при помощи цикла \texttt{\{\{ range . \}\}} перебираются все элементы строки и при помощи
\texttt{<td>} записываются по очереди в каждую ячейку строки.
